\begin{proposition}
\label{proposition-pushout-along-closed-immersion-and-integral}
\begin{reference}
\cite[Theorem 7.1 part iii]{Ferrand-Conducteur}
\end{reference}
In Situation \ref{situation-pushout-along-closed-immersion-and-integral}
the pushout $Y \amalg_Z X$ exists in the category of schemes. Picture
$$
\xymatrix{
Z \ar[r]_i \ar[d]_j & X \ar[d]^a \\
Y \ar[r]^-b & Y \amalg_Z X
}
$$
The diagram is a fibre square, the morphism $a$ is integral,
the morphism $b$ is a closed immersion, and
$$
\mathcal{O}_{Y \amalg_Z X} =
b_*\mathcal{O}_Y \times_{c_*\mathcal{O}_Z} a_*\mathcal{O}_X
$$
as sheaves of rings where $c = a \circ i = b \circ j$.
\end{proposition}

\begin{proof}
As a topological space we set $Y \amalg_Z X$ equal to the pushout of
the diagram in the category of topological spaces (Topology, Section
\ref{topology-section-colimits}). This is just the pushout
of the underlying sets (Topology, Lemma \ref{topology-lemma-colimits})
endowed with the quotient topology.
On $Y \amalg_Z X$ we have the maps of sheaves of rings
$$
b_*\mathcal{O}_Y \longrightarrow c_*\mathcal{O}_Z
\longleftarrow a_*\mathcal{O}_X
$$
and we can define
$$
\mathcal{O}_{Y \amalg_Z X} =
b_*\mathcal{O}_Y \times_{c_*\mathcal{O}_Z} a_*\mathcal{O}_X
$$
as the fibre product in the category of sheaves of rings. To prove that we
obtain a scheme we have to show that every point has an
affine open neighbourhood. This is clear for points not in the image of $c$
as the image of $c$ is a closed subset whose complement is
isomorphic as a ringed space to $(Y \setminus j(Z)) \amalg (X \setminus i(Z))$.

\medskip\noindent
A point in the image of $c$ corresponds to a unique $y \in Y$
in the image of $j$. By
Lemma \ref{lemma-prepare-pushout-along-closed-immersion-and-integral}
we find affine opens $U \subset X$ and $V \subset Y$ with
$y \in V$ and $i^{-1}(U) = j^{-1}(V)$.
Since the construction of the first paragraph is clearly compatible
with restriction to compatible open subschemes, to prove that it
produces a scheme we may assume $X$, $Y$, and $Z$ are affine.

\medskip\noindent
If $X = \Spec(A)$, $Y = \Spec(B)$, and $Z = \Spec(C)$ are affine, then
More on Algebra, Lemma \ref{more-algebra-lemma-points-of-fibre-product}
shows that $Y \amalg_Z X = \Spec(B \times_C A)$ as topological spaces.
To finish the proof that $Y \times_Z X$ is a scheme, it suffices to show
that on $\Spec(B \times_C A)$ the structure sheaf is the fibre product
of the pushforwards. This follows by applying
More on Algebra, Lemma \ref{more-algebra-lemma-diagram-localize}
to principal affine opens of $\Spec(B \times_C A)$.

\medskip\noindent
The discussion above shows the scheme $Y \amalg_Z X$
has an affine open covering $Y \amalg_Z X = \bigcup W_i$
such that $U_i = a^{-1}(W_i)$, $V_i = b^{-1}(W_i)$, and
$\Omega_i = c^{-1}(W_i)$ are affine open in $X$, $Y$, and $Z$.
Thus $a$ and $b$ are affine.
Moreover, if $A_i$, $B_i$, $C_i$ are the rings corresponding to
$U_i$, $V_i$, $\Omega_i$, then $A_i \to C_i$ is surjective and
$W_i$ corresponds to $A_i \times_{C_i} B_i$ which surjects onto
$B_i$. Hence $b$ is a closed immersion.
The ring map $A_i \times_{C_i} B_i \to A_i$ is integral by
More on Algebra, Lemma \ref{more-algebra-lemma-fibre-product-integral}
hence $a$ is integral. The diagram is cartesian because
$$
C_i \cong B_i \otimes_{B_i \times_{C_i} A_i} A_i
$$
This follows as $B_i \times_{C_i} A_i \to B_i$
and $A_i \to C_i$ are surjective maps whose kernels are the same.

\medskip\noindent
Finally, we can apply Lemmas \ref{lemma-basic-example-pushout} and
\ref{lemma-pushout-fpqc-local} to conclude our construction is a pushout
in the category of schemes.
\end{proof}